\chapter{Key Concepts for Developers}
\label{chap-key-concepts}

\renewcommand{\headrulewidth}{0.3pt}
\fancyhead[LE]{\thepage} \fancyhead[RE]{\nouppercase{\textbf{\leftmark}}}
\fancyhead[RO]{\thepage} \fancyhead[LO]{\nouppercase{\textbf{\rightmark}}}
\fancyfoot[L,C,R]{}

\setlength{\parindent}{0.35cm}

\minitoc
\newpage
% ================================================================

This chapter gathers a small set of \textbf{behavioural rules} that the code
follows silently --- decisions that are not obvious from reading a single
function, but that a new developer must know before touching data flowing
between \cawaqs~and \cwv. It is the natural home for ``why does the code do
\emph{that}?'' knowledge; the architectural tour and the concrete editing
recipes follow in \cref{chap-dev}.

\section{Coupling Observation Points to Mesh Cells}
\label{sec-obs-coupling}

An observation point (a gauging station, a piezometer, \ldots) must be attached
to exactly one mesh cell before its simulated values can be read. \cwv~decides
which cell to use in \textbf{two different ways}, and the choice is made
\emph{per observation}, not globally:

\begin{itemize}
    \item \textbf{By identifier (preferred).} If the observation layer carries an
    explicit cell identifier --- a configured \script{ID\_ABS} column in the
    observation's attribute table (\script{.dbf}) --- that identifier is used
    directly. This is the authoritative path: the coupling is exactly what the
    data declares, independent of geometry.

    \item \textbf{By geometry (fallback).} If no such identifier is provided, the
    code falls back to a \textbf{nearest-cell} search: it takes the observation's
    point geometry and finds the closest cell of the relevant mesh layer through
    a spatial-index nearest-neighbour lookup.
\end{itemize}

\alertwarning{The two paths are \textbf{not} equivalent. The geometric fallback
is a convenience for data that lacks explicit identifiers; it can attach a point
to the wrong cell near mesh boundaries or where cells are small. Whenever
authoritative results matter, provide \script{ID\_ABS} on the observations and
rely on the identifier path --- the nearest-cell search should be treated as a
last resort, not the default.}

\section{The HYD Compartment: \texorpdfstring{\script{ID\_ABS}}{ID\_ABS} vs.\
\texorpdfstring{\script{ID\_GIS}}{ID\_GIS}}
\label{sec-hyd-id-translation}

For every compartment \emph{except one}, a cell keeps the \textbf{same
identifier} on both sides of the tool: the ID used inside \cawaqs~is the ID used
inside \cwv. This lets most of the code move IDs across the \cawaqs~/~\cwv~seam
without any conversion.

The \textbf{HYD (hydrological/river) compartment is the sole exception.} Its
cells are identified differently on each side:

\begin{itemize}
    \item inside \cawaqs, a HYD cell is identified by its \textbf{\script{ID\_ABS}}
    (the absolute/internal \cawaqs~index);
    \item inside \cwv~(the GIS world), the same cell is identified by its
    \textbf{\script{ID\_GIS}}.
\end{itemize}

The two are bridged by a \textbf{HYD correspondence table} produced with the
simulation outputs, which pairs each \script{ID\_GIS} with its \script{ID\_ABS}.
The translation is applied \textbf{in both directions}, at the two ends of the
data path:

\begin{itemize}
    \item \textbf{On input} (reading a HYD observation / feeding \cwv): the
    incoming \cwv-side \script{ID\_GIS} is looked up in the correspondence table
    and \textbf{translated to \script{ID\_ABS}} so that \cawaqs~results can be
    addressed.
    \item \textbf{On output} (reporting HYD results back to the GIS): the reverse
    translation is applied, mapping \script{ID\_ABS} back to \script{ID\_GIS}.
\end{itemize}

\alertinfo{Practical consequence: if you are debugging a HYD result and the IDs
``don't line up'', check whether you are comparing an \script{ID\_ABS}
(\cawaqs~side) against an \script{ID\_GIS} (\cwv~side). For every other
compartment the identifiers coincide and no such lookup happens; only HYD carries
--- and requires --- this correspondence table.}
